\section{Received Traditions: Authority Without a Modern Fact Decree}

The Church often receives a Marian tradition through liturgy, papal teaching, a saint's life, an approved sacramental, a shrine, and generations of pastoral fruit rather than through one surviving decree declaring an apparition supernatural. These forms of reception are neither negligible nor interchangeable with an event judgment. The four dossiers below ask what the Church has positively received and leave the unlocated or ambiguous juridic step unnamed.

\subsection{Guadalupe: appearance, image, and a mestiza Church}

\dossieraxis{The Mexican tradition of Mary's appearances to the Indigenous Christian Juan Diego Cuauhtlatoatzin in December 1531, the sign of roses, the image on his tilma, and the building of a shrine at Tepeyac.}{Repeated papal reception, Juan Diego's canonization, and the extension of Guadalupe's feast throughout the American continent establish an exceptionally strong ecclesial reception. This audit does not claim a located modern decree adjudicating the original event.}{Papal and continental liturgical reception warrant theological treatment of the tradition; they do not settle every historical question about the dating and transmission of each narrative source or every proposed reading of the image.}

The received Nahuatl narrative places the encounter ten years after the fall of Tenochtitlan, amid military conquest, epidemic disease, disrupted communities, forced cultural translation, and an evangelization entangled with colonial power. In that narrative Juan Diego, a Nahua layman, hears Mary address him with familial tenderness and ask that a ``little sacred house'' be built where she may show her love and aid to all the inhabitants of the land. When the bishop asks for a sign, Castilian roses appear out of season and the image is disclosed on the tilma.

The extant printed Nahuatl narrative appears in Luis Laso de la Vega's 1649 \emph{Huei tlamahuiçoltica}; authorship, possible earlier composition, and transmission remain debated in modern scholarship. A serious Catholic account need not solve that debate by decree. The ecclesially received meaning is nevertheless unusually rich. The woman is recognizably the Mother of the true God and bears visual language intelligible across worlds. Catholic reception does not reduce her either to a Christianized pre-Christian deity or to a merely imposed European image. The image has been received as an inculturated proclamation: the Gospel enters a people, purifies and fulfills what is good, judges conquest's violence, and gives rise to a genuinely local Catholic subject.

The word \emph{mestiza}, used by modern Popes for the Guadalupan Virgin and for the Church she symbolizes, names more than mixed physical traits. At Pentecost the one Gospel does not erase languages; it creates communion through them. Guadalupe likewise resists both cultural domination and a syncretism that dissolves Christ. Mary is fruitful precisely because she is Jewish, Mother of the incarnate Word, and mother in the new people constituted by him. Her maternal face announces that Indigenous persons are not raw material for empire but hearers of God, agents of evangelization, and members of the Church.

The basilica's location and the bishop's role give the tradition an ecclesiology. Juan Diego's testimony must be heard; the pastor must discern; the sign is given for a house where a people may gather. Neither lay charism nor office is sufficient alone. The poor evangelize the hierarchy, and the hierarchy tests and receives the charism for the Church. John Paul II therefore called Guadalupe a model of perfectly inculturated evangelization and invoked her as ``Mother and Evangelizer of America'' and ``Star of the first and new evangelization.''

Guadalupe has also been invoked for human dignity, the unborn, migrants, peoples divided by race and class, and the unity of the Americas. These developments are legitimate when they arise from the Incarnation and protect concrete persons. They become distortions when the image is nationalized against another people, made party property, or used to romanticize the conquest. The Mother who forms one people in Christ cannot sponsor contempt.

The tilma is a sign, not a second Scripture or a puzzle requiring endlessly novel ``scientific'' discoveries. Claims about microscopic figures, impossible pigments, temperature, heartbeat, or astronomical codes vary greatly in evidential quality and lie outside the received core unless carefully verified. The deepest iconography is public: a mother bears Christ toward a wounded world, the lowly are heard, and a house of communion is built.

\subsection{Rue du Bac: the Miraculous Medal and grace without magic}

\dossieraxis{Reported encounters of St.~Catherine Labouré with the Blessed Virgin at the Daughters of Charity chapel on Rue du Bac, Paris, in 1830, especially the vision that supplied the design of the medal later called ``Miraculous.''}{Archbishop Hyacinthe-Louis de Quélen authorized production of the medal in 1832 after investigation; its diffusion, the 1842 Ratisbonne conversion, Catherine's canonization, liturgical reception, papal use, and universal devotion constitute authoritative reception. No conventional positive apparition decree is claimed here.}{The Church receives the medal and its orthodox Marian symbolism. She does not guarantee every later legend, automatic promise, or physical claim attached to it.}

The experiences came in a France still marked by the Revolution, the restoration and fall of monarchies, industrial poverty, anticlerical conflict, and the July Revolution of 1830. Catherine was a novice in a community founded to serve the poor. The setting matters: the image was entrusted not to an esoteric circle but within a vocation that joined contemplation to hospitals, foundlings, the elderly, and the destitute.

The medal's obverse shows Mary standing upon the globe and crushing the serpent, rays streaming from jeweled rings on her hands, surrounded by the prayer, ``O Mary, conceived without sin, pray for us who have recourse to thee.'' The reverse joins the letter M and a cross with the two Hearts and twelve stars. The whole is a compressed biblical catechesis: Genesis's enmity with the serpent, the woman and the dragon of Revelation, the Cross, the pierced hearts, the apostolic people, and Mary's intercession as a creature filled with grace.

The Immaculate Conception appears here twenty-four years before its solemn definition, but the medal did not create the doctrine. The Church had prayed, argued, and developed its understanding for centuries; the 1854 definition rests on public Revelation as received in Tradition. The devotion helped dispose the faithful to love the truth, while the dogma supplies the norm by which the devotion is interpreted. Private revelation serves doctrine; it does not promulgate it.

The rays signify graces obtained through Mary's intercession, yet their source is the Triune God through Christ. Rings from which no rays fall have been popularly read as graces not asked for; that interpretation can encourage prayer, but it cannot imply divine stinginess or make Mary's action independent. The reverse's M does not place Mary on the same plane as the Cross. It shows maternal union and discipleship beneath the one redemptive sacrifice.

Rapid diffusion and reported favors gave the medal its popular name. Alphonse Ratisbonne was wearing one when his 1842 conversion followed the silent apparition at Sant'Andrea delle Fratte, linking two dossiers without merging their judgments. ``Miraculous'' describes providential fruit; it does not turn the object into a talisman. A blessed medal is a sacramental: through the Church's prayer it disposes the faithful to receive grace and sanctify circumstances. Its use requires faith, conversion, prayer, and charity, never a mechanical transaction.

Rue du Bac's enduring synthesis is therefore active and ecclesial. The Immaculate Mother points to the Cross; received grace becomes service; a small reproducible object carries Christian memory into ordinary life. The most faithful wearer becomes available to the poor as Catherine did.

\subsection{Vailankanni: an Indian shrine of health and encounter}

\dossieraxis{Sixteenth-century traditions at Vailankanni in Tamil Nadu concerning Marian encounters with poor or sick boys, the carrying of milk, healing, and the growth of the shrine of Our Lady of Good Health; later tradition also associates Portuguese sailors' deliverance with the site.}{In 2024 the DDF praised the ``beautiful traditions,'' abundant fruits, interreligious attraction, and Christ-centered understanding of healing at the sanctuary. Its letter did not employ a formal result formula or declare the originating events supernatural.}{The received shrine tradition may be narrated as tradition and interpreted from the DDF's letter. Precise dates, conflated versions, and unverified miracle stories are not elevated into adjudicated facts.}

Vailankanni emerged on the Coromandel Coast in a world of Tamil villages, maritime exchange, Portuguese presence, illness, and unequal social encounters. The early narratives place humble children at the center. In one, a boy carrying milk meets a beautiful woman with a child; in another, a disabled boy receives healing and a commission. The sanctuary that developed became a major Catholic pilgrimage site attracting Christians, Hindus, Muslims, and people of other or no explicit faith.

The DDF's 2024 letter makes the Christological center explicit. The Virgin of Good Health carries Jesus; every blessing is received from him. Mary's maternity brings the sick to the divine physician, but she is not a parallel healer and the shrine is not a technology for obtaining cures. Health in Christian usage includes bodily care, reconciliation, freedom from sin, restored community, and final salvation. The uncured remain at the center of the Church's compassion.

Vailankanni's cultural form is not a defect to be stripped away. Barefoot pilgrimage, local music, flags, processions, offerings, and Tamil, Malayalam, Konkani, English, and other languages can embody a Catholic faith at home in India. Inculturation again requires two movements: the Gospel receives human signs and purifies whatever conflicts with worship of the Triune God. Practices must not foster caste exclusion, commercial exploitation, or ambiguity about the difference between veneration and adoration.

The interreligious presence is a pastoral gift and a theological responsibility. Hospitality must be real; no visitor should be reduced to a mission statistic. At the same time, Catholic identity is not concealed: Mary is honored because she is Mother of Jesus Christ, and the Eucharist is the sanctuary's summit. Shared human petitions for health can become a threshold of encounter without claiming that all religions make identical assertions.

The DDF's warm letter illustrates authoritative reception without juridic inflation. It recognizes the same Mother of Jesus beloved at other shrines and celebrates the Spirit's fruit among pilgrims. It does not provide missing archival certainty about the sixteenth century. Vailankanni can be fully honored when both truths are said.

\subsection{Cuapa: prayer and peace amid revolution}

\dossieraxis{Experiences reported by Bernardo Martínez, a sacristan at Cuapa, Nicaragua, beginning in April 1980, with Marian encounters reported between May and October and a message centered on prayer, the Rosary, peace, and reconciliation.}{In his preface of 13 November 1982, Bishop Pablo Antonio Vega wrote that as local bishop he was obliged to assure the authenticity of the facts in order to discern the message; no separate formal supernatural-origin decree was located.}{Cuapa is therefore treated as an ambiguous but serious positive diocesan reception, not as a documented \latin{constat}. Its message must not be recruited as a divine endorsement of a Nicaraguan faction.}

The events came one year after the Sandinista revolution overthrew the Somoza dictatorship and as hopes for liberation, ideological conflict, foreign intervention, and armed struggle were hardening into the Contra war. The Catholic community itself contained profound disagreement over revolution, class, authority, and the relation between faith and politics. In such a setting any claim from heaven becomes vulnerable to propaganda.

The reported Marian message asks for prayer, especially the Rosary, the practice of the five first Saturdays, peace, forgiveness, and the making of peace rather than merely requesting it. This last emphasis is morally demanding. Peace is not passivity before oppression, a pious substitute for justice, or victory for one's side. Catholic peace rests on truth, protects the vulnerable, seeks a just political order, rejects vengeance, and recognizes the enemy as a person capable of conversion.

The Rosary gives that work a Christological form. The mysteries remember a Savior born under empire, a kingdom announced to the poor, a Passion suffered without retaliatory hatred, and a Resurrection that defeats violence without denying its wounds. Repetition can disarm the imagination formed by slogans and fear. It does not supply policy details or exempt citizens from prudential judgment.

Martínez's ordinary ministry as sacristan and his reported reluctance fit a recurring pattern in the corpus, but narrative modesty is not proof. Bishop Vega's pastoral support deserves honest weight; the ambiguity of its juridic formula deserves equal honesty. Later summaries often convert the preface into a formal approval without showing the decree. This study declines that conversion.

Cuapa's most credible ecclesial appropriation is accordingly nonpartisan and exacting: Christians of opposed camps submit their consciences to Christ, repent of injustice committed by their own side, intercede for the dead and endangered, protect families, and labor for reconciliation joined to truth. If the apparition is used to baptize a geopolitical program, the call to peace has already been lost.

\subsection{The theological value of differentiated reception}

Guadalupe, Rue du Bac, Vailankanni, and Cuapa cannot be placed under one juridic caption. Guadalupe has exceptionally strong papal and continental liturgical reception; Rue du Bac, an approved sacramental and universal devotional reception; Vailankanni, an unusually affirmative DDF appraisal of a shrine tradition; Cuapa, a positive but formally ambiguous diocesan text. Naming these differences does not drain them of grace. It prevents devotion from depending on an inaccurate historical claim.

The traditions also widen the geographical and cultural field. A Nahua neophyte, a French Daughter of Charity, Tamil village children, and a Nicaraguan sacristan stand in histories of conquest, revolution, poverty, colonial encounter, and war. Mary is not a free-floating archetype above those histories. Because she is Mother of the incarnate Word, authentic devotion enters cultures, lifts the lowly, and forms communion while remaining capable of judging every culture---including Catholic subcultures---by the Gospel.
